% ============================================================================
% related_work.tex
% Related-work synthesis for the "digital mice" project.
%
% Scope of the digital-mice pipeline (for comparison throughout):
%   (a) instructed cognitive deficits (CoT suppression and prompt-induced
%       bounded rationality) layered onto frontier LLMs;
%   (b) calibration to human auction data via Li (2017) moment matching;
%   (c) SMAD-only evolutionary fitness over auction genotypes.
%
% This file is self-contained: it provides its own \begin{thebibliography}
% and can be \input{} into the main manuscript or compiled standalone with
% a minimal preamble.
% ============================================================================

\section*{Related work: digital mice in context}

\subsection*{Executive summary}

The closest antecedents to ``digital mice'' fall into six threads.
First, recent work explicitly reframes cognitive bias and bounded
rationality as \emph{instruments} for computational social science
rather than defects to mitigate (CoBRA; \citealp{cobra2025}), and a
parallel line uses LLMs as ``Homo silicus'' --- prompt-conditioned
implicit models of humans to be experimented on
\citep{horton2023silicus, manning2024automated}. Second, the
silicon-sampling tradition \citep{argyle2023silicon} and its
digital-twin descendants \citep{peng2025funhouse} build synthetic
populations through demographic persona conditioning, but typically
without manipulating reasoning capacity. Third, an emerging
auctions-and-games literature deploys LLMs as economic agents and
benchmarks their behavior against the canonical human experimental
record \citep{shah2025synthetic, zheng2025bounded, henning2025bubbles},
in some cases formalizing bounded rationality through QRE
\citep{pechon2026qre} or prompt-level equilibrium concepts
\citep{zhu2025llmnash}. Fourth, Li (2017)-style OSP theory
\citep{li2017osp, breitmoser2022obviousness} provides the diagnostic
vocabulary --- obvious dominance, contingent reasoning failures ---
that the digital-mice fitness function reuses. Fifth, multiple
empirical studies document that off-the-shelf LLMs are
\emph{too} rational, too homogeneous, and too fair relative to humans
\citep{henning2025bubbles, peng2025funhouse}, which is precisely the
gap a controlled-deficit pipeline is built to close. Sixth, and
methodologically closest, an AI-safety literature already studies how
to make a strong model \emph{play dumb}: sandbagging and
password-locking deliberately suppress a capable model's competence
\citep{vanderweij2024sandbagging, greenblatt2024password}, and
``model organisms'' induce minimal controlled failure modes as
experimental systems \citep{turner2025modelorganisms}. Digital mice
borrow this play-dumb move but repurpose it from a safety threat or
interpretability probe into a social-science instrument.

Digital mice are novel along two axes that no prior line covers
jointly: (i) systematic, prompt-induced reasoning-capacity knockouts
treated as an experimental design variable (a ``knockout mouse'' for
LLM agents), and (ii) quantitative calibration of those knockouts
against Li (2017) auction moments, then used as the fitness signal in
an evolutionary search over auction genotypes. Prior work has the
ingredients --- persona conditioning, bias toolkits, auction
benchmarks, QRE calibration --- but none combines instructed
cognitive deficits with moment-matched human auction data inside a
mechanism-design search loop.

% ----------------------------------------------------------------------------
\subsection*{1. Bias and bounded rationality as research instruments}

A small but growing literature reframes LLM cognitive limitations as
something to \emph{program in}, not engineer away. CoBRA
\citep{cobra2025} provides a toolkit with a Cognitive Bias Index and a
Behavioral Regulation Engine that explicitly controls biases in
LLM-based social agents, using validated social-science experiments
as the scaffold, and positions ``bias not as a flaw but as an
instrument for scientific reproducibility and behavioral modeling.''
This is the same instrumental move the digital-mice analogy makes:
treat the deficit as the apparatus, not the noise. The contrast is
along the deficit axis --- CoBRA targets classical cognitive biases
(anchoring, framing, etc.), whereas digital mice target
\emph{reasoning depth and contingent reasoning capacity} relevant to
auction strategy. Neither library calibrates against quantitative
auction moments.

\citet{zheng2025bounded} push in the same direction empirically:
across repeated games (RPS, Prisoner's Dilemma) LLMs reproduce
familiar human heuristics such as outcome-based strategy switching
and increased cooperation under repeated interaction, but ``apply
these rules more rigidly and demonstrate weaker sensitivity to the
dynamic changes in the game environment.'' They conclude that current
LLMs ``capture only a partial form of human-like bounded
rationality,'' motivating training methods for flexible opponent
modeling. Digital mice take the opposite engineering route: rather
than retraining for more flexible bounded rationality, we
\emph{prompt-instrument} specific deficits and treat the resulting
behavior as the unit of analysis.

% ----------------------------------------------------------------------------
\subsection*{2. Homo silicus and automated social science}

\citet{horton2023silicus} introduces the ``Homo silicus'' framing:
LLMs as ``implicit computational models of humans \dots\ they can be
given endowments, information, preferences, and so on, and then their
behavior can be explored in scenarios via simulation.'' Methodologically,
the paper replicates Charness--Rabin, Kahneman et al., Samuelson--Zeckhauser,
Oprea, and Horton himself, and treats \emph{qualitative} agreement (or
generative divergence) as the validation criterion. Digital mice keep
the prompt-as-endowment idea but replace qualitative replication with
quantitative moment matching against Li (2017).

\citet{manning2024automated} push this to the full research loop:
LLMs as both \emph{scientists} (hypothesis generators) and
\emph{subjects} (experimental agents), with structural causal models
serving as the language for stating hypotheses, building agents,
designing experiments, and analyzing data, demonstrated across
negotiations, bail hearings, job interviews, and auctions. Crucially,
the paper does not introduce controlled cognitive deficits or
bounded-rationality constraints --- agents are treated as default
reasoners. Digital mice can be read as a specialization of this
framework where the SCM ``treatment'' variable is reasoning capacity
itself, instrumented via prompted CoT suppression.

% ----------------------------------------------------------------------------
\subsection*{3. Synthetic populations and digital twins}

\citet{argyle2023silicon} coin ``silicon sampling'' and the criterion
of ``algorithmic fidelity'': conditioning GPT-3 on thousands of
sociodemographic backstories from real survey respondents, with the
goal that ``proper conditioning will cause it to accurately emulate
response distributions from a wide variety of human subgroups.'' The
construct propagates downstream into the digital-twin literature.
Methodologically, this is persona/demographic conditioning, with no
explicit treatment of bounded rationality or reasoning-capacity
ablation --- a clear contrast with the digital-mice approach, which
holds persona fixed and varies cognitive capacity.

\citet{peng2025funhouse} stress-test the digital-twin program across
nineteen pre-registered replications (164 outcomes) and report an
average correlation of only $r \approx 0.20$ with human responses, with
documented distortions including a ``Targeting Fairness'' regime in
which LLMs produce systematically skewed fairness judgments in
economic and social scenarios. The finding is directly relevant to
LLMs-as-bidders or LLMs-as-matching-market participants, where
fairness and envy considerations shape bid shading and preference
reports. Digital mice respond to this critique by refusing the naive
``proxy'' framing: rather than asking whether LLMs \emph{are} humans,
we ask whether instructed-deficit LLMs can reproduce the
\emph{moments} of human play under specific cognitive constraints.

% ----------------------------------------------------------------------------
\subsection*{4. LLMs as auction and game-theoretic agents}

\citet{shah2025synthetic} deploys LLMs as bidders in canonical
experimental auction frameworks (first-price, second-price, and OSP
variants), running over a thousand auctions and reporting that LLM
bidders ``produce results consistent with risk-averse human bidders''
and ``succumb to the winner's curse in common value settings.'' The
paper explicitly engages \citet{li2017osp} on obviously strategy-proof
mechanisms and \citet{breitmoser2022obviousness} on ``obviousness'' as
a theoretical lens for LLM cognitive limitations, and quantitatively
benchmarks against \citet{coppinger1980incentives},
\citet{cox1988theory}, and \citet{kagel1993independent} --- e.g.,
matching ``the 27\% reported by Kagel (1993)'' for value-bidding
frequency in SPSB. This is the most direct precedent for the
auction-side empirical loop in digital mice, sharing both the OSP
diagnostic language and the human-baseline calibration target. The
key difference: \citet{shah2025synthetic} does \emph{not} systematically
manipulate reasoning capacity to recover specific moments, and does not
run an evolutionary search over auction formats conditional on agent
cognition.

\citet{pechon2026qre} provide a complementary calibration tool by
estimating Quantal Response Equilibrium $\lambda$ parameters for LLM
agents on a continuous scale against human data (with human
$\lambda \in [1.0, 2.5]$ and frontier-model estimates in
$[0.05, 1.10]$ across 1{,}855 games and seven models). QRE gives a
principled one-dimensional handle on strategic sophistication, and is
in principle complementary to digital mice: QRE measures
\emph{how} rational an LLM is, while digital mice manipulate that
quantity through instructed deficits.

\citet{zhu2025llmnash} formalizes an ``LLM-Nash reasoning equilibrium''
where agents choose \emph{reasoning prompts} rather than actions, and
defines mindsets as triples (information intake, prompt space,
worldview parameters), with equilibrium defined over the prompt space
and actions emerging as LLM inference output. This provides a clean
theoretical vocabulary for what digital mice do mechanically:
deficit-injecting prompts are precisely the strategic objects in
$\mathcal{X}, \mathcal{Y}$ above, and the equilibrium concept formalizes
why prompt-level manipulations produce behavior distinct from
classical Nash play.

% ----------------------------------------------------------------------------
\subsection*{5. Out-of-the-box LLMs are too rational and too homogeneous}

A robust empirical finding across studies is that frontier LLMs,
absent intervention, are systematically \emph{more} rational and
\emph{more} homogeneous than human subjects.
\citet{henning2025bubbles} run the Smith--Suchanek--Williams bubble
paradigm with LLM populations and find that LLMs ``generally exhibit
a \emph{textbook-rational} approach, pricing the asset near its
fundamental value, and show only a muted tendency toward bubble
formation,'' and that LLM-based agents ``display less trading strategy
variance in contrast to humans.'' They warn that ``LLM-only data
\dots\ to replicate human-driven market phenomena'' is unreliable
because ``key behavioral features, such as large emergent bubbles,
were not robustly reproduced.'' Together with the fairness
distortions documented by \citet{peng2025funhouse} and the rigidity
findings of \citet{zheng2025bounded}, this body of evidence motivates
the central digital-mice design choice: do not assume a baseline LLM
matches a human bidder; instead, calibrate a \emph{cognitive
treatment} that reproduces specified moments of human play, and use
that treated population as the unit of analysis.

% ----------------------------------------------------------------------------
\subsection*{6. Researcher-controlled deficit induction: making a strong
model ``play dumb'' on purpose}

The thread closest to the digital-mice \emph{mechanism}---taking a
capable reasoning model and deliberately degrading its competence in a
controlled way---divides cleanly into two classes that should not be
conflated. \emph{(A) Model-driven evasion}: the model itself, or a
developer with an incentive to pass an evaluation, hides capability.
\emph{(B) Researcher-driven controlled induction}: an experimenter
systematically dials a model's competence down as a designed treatment,
to study or use the resulting deficient agent. Digital mice are squarely
in class (B). We review (A) only to delimit it, then survey the (B)
methods that are the true methodological neighbors.

\paragraph{Delimiting the evasion case (A).}
Sandbagging is the canonical (A) example: \citet{vanderweij2024sandbagging}
define it as ``strategic underperformance on an evaluation'' and show
frontier models can be prompted to ``selectively underperform on
dangerous capability evaluations'' while passing general ones, and that
``capable password-locked models can reasonably emulate less capable
models.'' \citet{greenblatt2024password} train password-locked models
that ``imitate a much weaker LLM'' unless a password is present. These
papers share the digital-mice \emph{technique}---a strong model made to
act weak---but their \emph{object of study} is the threat that a model
or developer evades oversight; the suppression is in service of
deception, not measurement. We borrow the one transferable empirical
lesson---suppressed capability is fragile and can leak back ``from a few
high-quality demonstrations''---because it is exactly what the
digital-mice reasoning-chain audit quantifies (the $0.001$ vs.\ $0.41$
dominance-invocation rate under CoT vs.\ instruction-only suppression).

\paragraph{Controlled depth limitation by prompt (B).}
The closest (B) method to the digital-mice contingent/belief axes is
level-$k$ and cognitive-hierarchy prompting. \citet{zhang2024klevel}
build a ``K-Level Reasoning'' framework in which the experimenter sets
the strategic depth---reason as a level-$0$, level-$1$, \dots, level-$k$
player forming ``higher order beliefs, beliefs about others' beliefs''---
and the cognitive-hierarchy benchmark \citep{chbench2025} evaluates LLMs
across that explicit depth ladder. This is researcher-controlled
bounded rationality of exactly the kind digital mice impose, drawn from
the same level-$k$ / cognitive-hierarchy tradition. The difference is
direction of use: K-level reasoning raises depth to \emph{improve}
strategic play, whereas digital mice hold depth low to \emph{reproduce}
a human deficit and then calibrate it to auction moments.

\paragraph{Controlled compute/reasoning-budget limitation (B).}
A second (B) method limits how much a model is allowed to reason. The
test-time-compute literature \citep{reasoningbudget2025survey} controls
a model's reasoning under a fixed budget---e.g.\ BudgetThinker
\citep{budgetthinker2025} inserts control tokens that tell the model its
remaining token budget, and production systems expose a ``thinking
budget'' knob. Cutting the budget is a continuous, researcher-set
deficit dial. Digital mice instead constrain \emph{which reasoning
operations} are permitted (contingent, forward, belief) rather than
\emph{how many tokens} are spent, but both treat reasoning capacity as a
controllable experimental quantity.

\paragraph{Controlled removal at the weight/representation level (B).}
Where digital mice ablate by prompt, machine unlearning ablates by
weights. RMU \citep{wmdp2024} ``reduces model performance'' on a target
capability ``while maintaining general capabilities,'' by redirecting
the internal representation of the forgetting set---a surgical,
researcher-controlled capability knockout. Activation steering and
persona vectors \citep{chen2025personavectors} achieve a similar
controlled effect at inference time: behaviors and traits correspond to
``approximately linear directions in activation space,'' and adding or
\emph{subtracting} along a direction reliably steers (or suppresses)
behavior ``without changing the surface prompt.'' These are the
weight- and activation-level analogues of the digital-mice
prompt-level deficit; they show the controlled-deficit philosophy is
not specific to prompting.

\paragraph{Model organisms: the exact framing twin (B).}
Finally, the alignment literature's \emph{model organisms} are the
precise methodological twin of the knockout-mouse framing:
deliberately induced, minimal, reproducible failure modes used as
controlled experimental systems. \citet{turner2025modelorganisms}
create model organisms of emergent misalignment using ``a single
rank-1 LoRA adapter,'' isolating ``a single linear direction that
induces'' the behavior so defensive techniques can be benchmarked
against a known ground truth. The shared logic is exact: induce a
controlled deficit, verify it is the intended one, use the resulting
strain to stress-test something downstream. Digital mice differ only
in the downstream target (a mechanism's SMAD against human auction
data, not an interpretability method) and the induction channel
(prompt-level CoT constraints, not weight-level edits).

\paragraph{What is distinctive about the digital-mice variant of (B).}
Across the (B) methods above, the deficit is induced for a
\emph{model-facing} purpose: better strategic play (level-$k$),
cheaper inference (budget control), safety (unlearning), or
interpretability (model organisms). None calibrates the induced
deficit to a \emph{human behavioral target}. Digital mice are the only
(B) instance we found that (i) induces the deficit specifically to
match human economic-experiment moments and (ii) then uses the
calibrated deficient population as the evaluator in a mechanism search.

% ----------------------------------------------------------------------------
\subsection*{7. Where digital mice are novel vs.\ reuse}

\paragraph{Reused ingredients.}
The instrumental framing of bias/bounded rationality
\citep{cobra2025}; the Homo-silicus prompt-as-endowment paradigm
\citep{horton2023silicus, manning2024automated}; the LLM-as-auction-bidder
setup and OSP diagnostic vocabulary
\citep{shah2025synthetic, li2017osp, breitmoser2022obviousness}; the
human-benchmark targets from classical experimental auctions
\citep{coppinger1980incentives, cox1988theory, kagel1993independent};
the quantitative measurement of LLM bounded rationality via QRE
\citep{pechon2026qre} or prompt-space equilibrium \citep{zhu2025llmnash};
and---most directly---the practice of making a capable model emulate a
weaker one, established as sandbagging / password-locking
\citep{vanderweij2024sandbagging, greenblatt2024password} and as
controlled deficit induction in model organisms
\citep{turner2025modelorganisms}.

\paragraph{Novel combination.}
Digital mice (i) treat reasoning-capacity manipulation (CoT
suppression, instructed deficits) as a designed experimental variable
in the spirit of a knockout mouse---the same ``play-dumb'' move as
sandbagging \citep{vanderweij2024sandbagging} and model organisms
\citep{turner2025modelorganisms}, but turned from a safety \emph{threat}
or interpretability \emph{probe} into a social-science \emph{instrument};
(ii) calibrate the induced-deficit population against Li (2017) auction
moments rather than against qualitative replication
\citep{horton2023silicus}, single-dimensional QRE \citep{pechon2026qre},
or a generic weaker-model target \citep{greenblatt2024password}; and
(iii) plug the calibrated population into an evolutionary search over
auction genotypes with an SMAD-only fitness function --- a
mechanism-design loop with no antecedent in any of the threads above.
No prior line combines deliberate capability suppression, calibration to
human economic-experiment moments, and mechanism search; that triple is
what is new.

% ============================================================================
\begin{thebibliography}{99}

\bibitem[Argyle et al.(2023)]{argyle2023silicon}
Argyle, L.~P., Busby, E.~C., Fulda, N., Gubler, J.~R., Rytting, C.,
and Wingate, D. (2023).
Out of one, many: Using language models to simulate human samples.
\emph{Political Analysis}, 31(3):337--351. arXiv:2209.06899.

\bibitem[Breitmoser and Schweighofer-Kodritsch(2022)]{breitmoser2022obviousness}
Breitmoser, Y. and Schweighofer-Kodritsch, S. (2022).
Obviousness around the clock.
\emph{Experimental Economics}, 25:1--29.

\bibitem[CoBRA(2025)]{cobra2025}
Authors of CoBRA (2025).
CoBRA: A toolkit for specifying and controlling cognitive biases in
LLM-based social agents.
\emph{Proceedings of CHI 2026} (Best Paper). arXiv:2509.13588.

\bibitem[Coppinger et al.(1980)]{coppinger1980incentives}
Coppinger, V.~M., Smith, V.~L., and Titus, J.~A. (1980).
Incentives and behavior in English, Dutch, and sealed-bid auctions.
\emph{Economic Inquiry}, 18(1):1--22.

\bibitem[Cox et al.(1988)]{cox1988theory}
Cox, J.~C., Smith, V.~L., and Walker, J.~M. (1988).
Theory and individual behavior of first-price auctions.
\emph{Journal of Risk and Uncertainty}, 1(1):61--99.

\bibitem[Greenblatt et al.(2024)]{greenblatt2024password}
Greenblatt, R., Roger, F., Krasheninnikov, D., and Krueger, D. (2024).
Stress-testing capability elicitation with password-locked models.
\emph{Advances in Neural Information Processing Systems (NeurIPS)}.
arXiv:2405.19550.

\bibitem[Henning et al.(2025)]{henning2025bubbles}
Henning, J., Ojha, A., Spoon, R., Han, J., and Camerer, C.~F. (2025).
LLM agents do not replicate human market traders: Evidence from
experimental finance.
arXiv:2502.15800.

\bibitem[Liu et al.(2025)]{liu2025cursecot}
Liu, T., et al. (2025).
The curse of CoT: On the limitations of chain-of-thought in in-context
learning.
arXiv:2504.05081.

\bibitem[Horton(2023)]{horton2023silicus}
Horton, J.~J. (2023).
Large language models as simulated economic agents: What can we
learn from \emph{Homo silicus}?
NBER Working Paper 31122. arXiv:2301.07543.

\bibitem[Kagel(1993)]{kagel1993independent}
Kagel, J.~H. (1993).
Independent private value auctions: Bidder behavior in first-,
second- and third-price auctions with varying numbers of bidders.
\emph{Economic Journal}, 103(419):868--879.

\bibitem[Kagel and Levin(1986)]{kagel1986winner}
Kagel, J.~H. and Levin, D. (1986).
The winner's curse and public information in common value auctions.
\emph{American Economic Review}, 76(5):894--920.

\bibitem[Li(2017)]{li2017osp}
Li, S. (2017).
Obviously strategy-proof mechanisms.
\emph{American Economic Review}, 107(11):3257--3287.

\bibitem[Manning et al.(2024)]{manning2024automated}
Manning, B.~S., Zhu, K., and Horton, J.~J. (2024).
Automated social science: Language models as scientist and subjects.
arXiv:2404.11794.

\bibitem[Pechon-Elkins and Chun(2026)]{pechon2026qre}
Pechon-Elkins, L. and Chun, S. (2026).
Quantal response equilibrium as a measure of strategic sophistication:
Theory and validation for LLM evaluation.
arXiv:2603.10029.

\bibitem[Peng et al.(2025)]{peng2025funhouse}
Peng, A. et al. (2025).
Digital twins as funhouse mirrors: A pre-registered evaluation of
LLM-based simulated participants.
arXiv:2509.19088.

\bibitem[Shah et al.(2025)]{shah2025synthetic}
Shah, A., Zhu, K., Jiang, Y., Wang, J.~G., Dayi, A.~K., Horton, J.~J.,
and Parkes, D.~C. (2025).
Learning from synthetic labs: Language models as auction participants.
arXiv:2507.09083.

\bibitem[Turner et al.(2025)]{turner2025modelorganisms}
Turner, E., Soligo, A., Taylor, M., Rajamanoharan, S., and Nanda, N.
(2025).
Model organisms for emergent misalignment.
arXiv:2506.11613.

\bibitem[van der Weij et al.(2024)]{vanderweij2024sandbagging}
van der Weij, T., Hofst\"atter, F., Jaffe, O., Brown, S.~F., and
Ward, F.~R. (2024).
AI sandbagging: Language models can strategically underperform on
evaluations.
arXiv:2406.07358.

\bibitem[Zheng et al.(2025)]{zheng2025bounded}
Zheng, X., Zhou, J., and Wang, Y. (2025).
Beyond Nash equilibrium: Bounded rationality of LLMs and humans in
strategic decision-making.
arXiv:2506.09390.

\bibitem[Zhu(2025)]{zhu2025llmnash}
Zhu, Q. (2025).
LLM-Nash equilibrium: Reasoning prompts as strategies.
arXiv:2507.08208.

% Researcher-controlled deficit-induction methods (Section 6, class B).

\bibitem[Zhang et al.(2024)]{zhang2024klevel}
Zhang, Y., Mao, S., Ge, T., Wang, X., Xia, Y., Lan, M., and Wei, F.
(2024).
K-level reasoning: Establishing higher order beliefs in large language
models for strategic reasoning.
arXiv:2402.01521.

\bibitem[CHBench(2025)]{chbench2025}
Authors of CHBench (2025).
CHBench: A cognitive hierarchy benchmark for evaluating strategic
reasoning capability of LLMs.
arXiv:2508.11944.

\bibitem[Reasoning-on-a-Budget survey(2025)]{reasoningbudget2025survey}
Authors of the test-time-compute survey (2025).
Reasoning on a budget: A survey of adaptive and controllable test-time
compute in LLMs.
arXiv:2507.02076.

\bibitem[BudgetThinker(2025)]{budgetthinker2025}
Authors of BudgetThinker (2025).
BudgetThinker: Empowering budget-aware LLM reasoning with control
tokens.
arXiv:2508.17196.

\bibitem[Li et al.(2024)]{wmdp2024}
Li, N., Pan, A., et al. (2024).
The WMDP benchmark: Measuring and reducing malicious use with
unlearning.
\emph{Proceedings of ICML}. arXiv:2403.03218.

\bibitem[Chen et al.(2025)]{chen2025personavectors}
Chen, R., et al. (2025).
Persona vectors: Monitoring and controlling character traits in
language models.
arXiv:2507.21509.

% Supporting / contextual references used in the synthesis but not yet
% expanded above; kept in the bibliography for completeness.

\bibitem[Charness and Rabin(2002)]{charness2002understanding}
Charness, G. and Rabin, M. (2002).
Understanding social preferences with simple tests.
\emph{Quarterly Journal of Economics}, 117(3):817--869.

\bibitem[Kahneman et al.(1986)]{kahneman1986fairness}
Kahneman, D., Knetsch, J.~L., and Thaler, R.~H. (1986).
Fairness as a constraint on profit seeking: Entitlements in the market.
\emph{American Economic Review}, 76(4):728--741.

\bibitem[Samuelson and Zeckhauser(1988)]{samuelson1988status}
Samuelson, W. and Zeckhauser, R. (1988).
Status quo bias in decision making.
\emph{Journal of Risk and Uncertainty}, 1(1):7--59.

\bibitem[Oprea(2024)]{oprea2024decisions}
Oprea, R. (2024).
Decisions under risk are decisions under complexity.
\emph{American Economic Review}, forthcoming.

\bibitem[Coppinger et al.(1980)]{coppinger1980supp}
See \citet{coppinger1980incentives}.

\end{thebibliography}
